\section*{Arquitectura -- Diagrama de bloques funcional}

En la Figura \ref{fig:bloques_hw}, los bloques en línea continua representan el hardware de la \textbf{Fase 1 (Semestre Actual)}, mientras que los bloques en naranja con línea punteada corresponden a la \textbf{Fase 2 (Tesis)}.

\begin{figure}[H]
\centering
\begin{tikzpicture}[node distance=1.2cm and 1.3cm, auto, font=\sffamily\small,
    mcu/.style={rectangle, draw=unalnavy, fill=blue!8, thick, minimum width=3.2cm, minimum height=3.6cm, align=center, rounded corners},
    block1/.style={rectangle, draw=black, fill=phase1blue, thick, minimum width=2.6cm, minimum height=1.1cm, align=center, rounded corners},
    block2/.style={rectangle, draw=unalred, dashed, fill=phase2orange, thick, minimum width=2.6cm, minimum height=1.1cm, align=center, rounded corners}]

    % MCU CENTRAL
    \node [mcu] (esp) {\textbf{NÚCLEO CENTRAL}\\[0.1cm]\textbf{ESP32-S3-N16R8}\\(Dual-Core 240 MHz)\\[0.1cm]\emph{3.3V Logic}};

    % ENTRADAS
    \node [block1, left=of esp, yshift=1.2cm] (gps) {GPS MAX-M10S\\(UART1: GPIO 16/17)};
    \node [block1, left=of esp, yshift=-1.2cm] (imu) {IMU LSM6DSOX\\(I$^2$C: GPIO 21/22)};
    
    % ENERGIA
    \node [block1, below=of esp, yshift=-0.2cm] (pwr) {MPPT CN3791 6V\\+ LiPo 3.7V + MAX17048};

    % SALIDAS Y TELEMETRIA
    \node [block1, right=of esp, yshift=1.4cm] (lora) {LoRa SX1262 (915 MHz)\\(SPI: GPIO 5/18/19/23)};
    \node [block1, right=of esp, yshift=0.1cm] (opto) {VO617A Optoacoplador\\+ Buzzer / LED (GPIO 13)};
    
    % EXPANSION FASE 2
    \node [block2, right=of esp, yshift=-1.3cm] (gsm) {Módem Quectel BG95-M3\\(UART2: GPIO 26/27)\\[0.05cm]\footnotesize\textbf{[Fase 2 -- Tesis]}};

    % CONEXIONES
    \draw[<->, thick] (gps) -- (esp);
    \draw[<->, thick] (imu) -- (esp);
    \draw[<->, thick] (esp) -- (lora);
    \draw[->, thick] (esp) -- (opto);
    \draw[->, thick] (pwr) -- (esp);
    \draw[<->, thick, dashed, draw=unalred] (esp) -- (gsm);

\end{tikzpicture}
\caption{Diagrama de bloques funcional de la arquitectura de hardware DR-VF-01.}
\label{fig:bloques_hw}
\end{figure}

\subsection*{Mapeo de Pines y Asignación de Periféricos}

\begin{table}[H]
\centering
\small
\begin{tabularx}{\textwidth}{|l l l X|}
\hline
\rowcolor{categorybg}
\textbf{Módulo / Periférico} & \textbf{Protocolo} & \textbf{Pines ESP32-S3} & \textbf{Notas de Ruteo en PCB} \\ \hline
GPS u-blox MAX-M10S & UART1 & GPIO 16 (RX1), GPIO 17 (TX1) & Intercambiar TX/RX en ruteo de pista. \\ \hline
IMU ST LSM6DSOX & I$^2$C & GPIO 21 (SDA), GPIO 22 (SCL) & Requiere resistencias pull-up de $4.7~\text{k}\Omega$ a 3.3V. \\ \hline
Medidor MAX17048 & I$^2$C & GPIO 21 (SDA), GPIO 22 (SCL) & Comparte bus I$^2$C con la IMU. \\ \hline
LoRa SX1262 & SPI & GPIO 5 (CS), 18 (SCK), 19 (MISO), 23 (MOSI) & Líneas DIO1 (GPIO 4) y BUSY (GPIO 2) mapeadas. \\ \hline
Optoacoplador / Buzzer & GPIO Digital & GPIO 13 (Disparo), GPIO 14 (PWM) & Separación estricta de plano de masa. \\ \hline
\rowcolor{phase2orange}
Quectel BG95-M3 [Fase 2] & UART2 & GPIO 26 (RX2), GPIO 27 (TX2) & \textbf{Pines reservados en PCB para expansión de tesis.} \\ \hline
\end{tabularx}
\end{table}

\subsection*{Árbol de Distribución de Potencia y Gestión Energética}

\begin{figure}[H]
\centering
\begin{tikzpicture}[node distance=0.8cm and 0.8cm, auto, font=\sffamily\small,
    pwr_in/.style={rectangle, draw=black, fill=categorybg, thick, text width=2.4cm, align=center, minimum height=1.1cm, rounded corners},
    pwr_mgt/.style={rectangle, draw=unalnavy, fill=phase1blue, thick, text width=2.6cm, align=center, minimum height=1.1cm, rounded corners},
    pwr_bus/.style={rectangle, draw=black, fill=blue!12, thick, text width=2.6cm, align=center, minimum height=1.1cm, rounded corners},
    pwr_load/.style={rectangle, draw=black, fill=white, thick, text width=2.6cm, align=center, minimum height=1.2cm, rounded corners},
    pwr_fase2/.style={rectangle, draw=unalred, dashed, fill=phase2orange, thick, text width=2.6cm, align=center, minimum height=1.1cm, rounded corners}]

    % NODO CENTRAL DE BATERÍA Y ENTRADAS (FILA 1)
    \node [pwr_bus] (bat) {Batería LiPo 3.7V\\(1200--1500 mAh)};
    \node [pwr_mgt, left=0.8cm of bat] (cn) {Cargador MPPT\\CN3791 (6V)};
    \node [pwr_in, left=0.8cm of cn] (in) {Entrada 6V Solar /\\ USB-C (5V)};

    % REGULADORES Y MÓDEM FASE 2 (FILA 2)
    \node [pwr_mgt, below=1.5cm of bat, xshift=-2.6cm] (ldo) {Regulador LDO\\3.3V (500 mA)};
    \node [pwr_fase2, below=1.5cm of bat, xshift=2.6cm] (bg95) {Quectel BG95-M3\\(Picos 400--600 mA)\\[0.05cm]\footnotesize\textbf{[Fase 2 -- Tesis]}};

    % CARGAS DE 3.3V ALINEADAS DEBAJO DEL LDO (FILA 3)
    \node [pwr_load, below=1.4cm of ldo] (sens) {Sensores \& Radio\\(MAX-M10S, LSM6DSOX\\SX1262, MAX17048)};
    \node [pwr_load, left=0.5cm of sens] (mcu) {ESP32-S3\\(Microcontrolador)};
    \node [pwr_load, right=0.5cm of sens] (act) {Disuasión Aislada\\(Opto VO617A / Buzzer)};

    % CONEXIONES DE ENTRADA Y CARGA
    \draw[->, thick] (in) -- (cn);
    \draw[->, thick] (cn) -- node[above, font=\footnotesize] {$V_{\text{CHG}}$} (bat);

    % DISTRIBUCIÓN DESDE BATERÍA A REGULADORES (FILA 1 -> FILA 2)
    % CORRECCIÓN: El texto V_BAT se ubica a la derecha de la línea vertical que baja
    \draw[thick] (bat.south) -- node[right, font=\footnotesize] {$V_{\text{BAT}}$ ($3.7\text{--}4.2\text{V}$)} ++(0,-0.8) coordinate (vbat_bus);
    \draw[->, thick] (vbat_bus) -| (ldo.north);
    \draw[->, thick, dashed, draw=unalred] (vbat_bus) -| (bg95.north);
    
    % BUS DE 3.3V DESDE LDO A CARGAS (FILA 2 -> FILA 3)
    \draw[thick] (ldo.south) -- ++(0,-0.5) coordinate (bus3v);
    \draw[->, thick] (bus3v) -| (mcu.north);
    \draw[->, thick] (bus3v) -- (sens.north);
    \draw[->, thick] (bus3v) -| (act.north);

    % Etiqueta 3.3V
    \node[above right, font=\footnotesize] at (bus3v) {$3.3\text{V}$};

\end{tikzpicture}
\caption{Árbol de distribución de potencia y gestión energética para el sistema DR-VF-01.}
\label{fig:arbol_potencia}
\end{figure}